
\subsubsection{Experimental Design}

We conduct systematic experiments across four log-spaced dataset scales (10K, 100K, 1M, 10M tokens) with five curation strategies to isolate individual effects and test compositional interactions.

\textbf{Dataset}: We sample from C4 realnewslike, a standard language model pretraining corpus derived from Common Crawl. Single-source sampling controls for domain shift confounds.

\textbf{Curation Strategies}:
\begin{enumerate}
\item Baseline: Random 49\% sampling (size control)
\item Q-only: V-Information filtering to 70\%, then random sample to 49\%
\item D-only: FAC-based diverse sampling to 70\%, then random sample to 49\%
\item QD (Quality-First): V-Information to 70\%, then FAC on that subset to 49\%
\item DQ (Diversity-First): FAC to 70\%, then V-Information on that subset to 49\%
\end{enumerate}

Fixed reduction rates (70\% first step, 49\% final) control for dataset size confounds---all strategies output the same token count, isolating ordering effect from reduction magnitude.

\subsubsection{Metrics}

\textbf{V-Information (Quality)}: Measures pointwise information content as V-Info(x) = -log p\_model(x) using a reference GPT-2 model. Documents above the 70th percentile are retained.

\textbf{Feature Activation Coverage (FAC, Diversity)}: Measures feature space coverage using intermediate layer activations. We select documents that maximize incremental FAC gain via greedy coverage optimization.

\textbf{Metric Orthogonality}: Correlation between V-Info and FAC scores on held-out 100K C4 sample: $\rho$=0.23 (p=0.08, marginally non-significant at $\alpha =0.05)$. This low correlation suggests metrics capture largely distinct aspects of data quality, though the marginal p-value indicates some shared variance cannot be ruled out.

\subsubsection{Model Training}

We use GPT-2 Small (124M parameters) as our proxy model. Data Mixing Laws validated that small proxy models predict large-scale mixture performance with high correlation, enabling computational feasibility.

\textbf{Training Configuration}:
\begin{itemize}
\item Epochs: 5 (validated as sufficient for convergence)
\item Optimizer: AdamW with learning rate 5e-4
\item Batch size: 32
\item Context length: 1024 tokens
\end{itemize}

\subsubsection{Evaluation}

We evaluate on three diverse downstream tasks:
\begin{enumerate}
\item MMLU (Massive Multitask Language Understanding): 5-shot accuracy
\item BEIR (Benchmark for Information Retrieval): nDCG@10, 0-shot
\item GSM8K (Grade School Math): exact match accuracy, 0-shot
\end{enumerate}

\textbf{Primary metric}: Average performance across all three benchmarks. We train 3 independent runs per condition with different random seeds and report mean $\pm$ standard deviation.

\subsubsection{Hypothesis Testing}

\textbf{H-E1 (Quality Saturation)}: Q-only improvement over baseline decreases monotonically with scale. Test via linear regression of $\Delta Q$ vs log(scale). Success: negative slope, p < 0.05.

\textbf{H-E2 (Diversity Persistence)}: D-only improvement increases with scale. Modeled from FAC literature showing $\rho$=0.90 correlation at large scale. Simulated trajectory pending full implementation due to environment constraints.

\textbf{H-M1 (Compositional Ordering)}: Sequential ordering (QD vs DQ) produces statistically different performance. Test via Cohen's d for (QD $-$ DQ) at all scales. Success: d > 0.5.

\textbf{H-C1 (Coefficient Trajectories)}: Quality coefficient $\alpha _Q$ decreases, diversity coefficient $\alpha _D$ increases with scale. Test via regression model Performance $\approx \alpha _{Q}\cdot Q + \alpha _{D}\cdot D$. Success: $\alpha _Q$ slope negative (p<0.05), $\alpha _D$ slope positive (p<0.05).
